\documentclass[11pt]{article}
\usepackage[letterpaper,margin=0.68in]{geometry}
\usepackage{amsmath,amssymb}
\usepackage{enumitem}
\usepackage{xcolor}
\usepackage{fancyhdr}
\usepackage{microtype}
\usepackage{tabularx,array,booktabs}
\usepackage{tikz}
\usepackage[most]{tcolorbox}

\definecolor{olive}{HTML}{4D5430}
\definecolor{gold}{HTML}{9A7B22}
\definecolor{pale}{HTML}{F5F1DC}
\definecolor{softgreen}{HTML}{EDF0E2}
\definecolor{softgray}{HTML}{F5F5F2}
\definecolor{ink}{HTML}{292D20}

\pagestyle{fancy}
\fancyhf{}
\lhead{\textcolor{olive}{MATH\& 107: Math in Society}}
\rhead{\textcolor{gold}{Fall 2026}}
\cfoot{\thepage}
\setlength{\headheight}{14pt}
\setlength{\parindent}{0pt}
\setlength{\parskip}{5pt}
\setlist[enumerate]{leftmargin=*,itemsep=7pt,topsep=5pt}
\setlist[itemize]{leftmargin=1.35em,itemsep=3pt,topsep=3pt}

\newtcolorbox{keybox}[1]{colback=pale,colframe=olive,boxrule=0.8pt,arc=2mm,title=\textbf{#1},fonttitle=\color{olive},left=2mm,right=2mm,top=1.5mm,bottom=1.5mm}
\newtcolorbox{examplebox}[1]{colback=softgreen,colframe=gold,boxrule=0.8pt,arc=2mm,title=\textbf{#1},fonttitle=\color{ink},left=2mm,right=2mm,top=1.5mm,bottom=1.5mm}
\newtcolorbox{refbox}[1]{colback=softgray,colframe=gold,boxrule=0.7pt,arc=2mm,title=\textbf{#1},fonttitle=\color{ink},left=2mm,right=2mm,top=1.5mm,bottom=1.5mm}

\newcommand{\summaryhead}[2]{%
  \clearpage
  {\Large\bfseries\textcolor{olive}{Module #1: #2 — Summary}}\par
  \vspace{2pt}\textcolor{gold}{\rule{\linewidth}{1.2pt}}\par\smallskip
}
\newcommand{\prequizhead}[2]{%
  \clearpage
  {\Large\bfseries\textcolor{olive}{Module #1: #2 — Prequiz}}\par
  \textit{Open-note. Show your mathematical work. A calculation and a clearly stated final answer are enough unless a sentence is specifically requested.}\par
  Name: \hspace{2.6in} Date: \hspace{1.1in}\par
  \vspace{2pt}\textcolor{gold}{\rule{\linewidth}{1.2pt}}\par\smallskip
}
\newcommand{\quizhead}[2]{%
  \clearpage
  {\Large\bfseries\textcolor{olive}{Module #1: #2 — Matching Quiz}}\par
  \textit{These problems use the same skills as the prequiz, with new values and contexts. Show enough work to make your reasoning visible.}\par
  Name: \hspace{2.6in} Date: \hspace{1.1in}\par
  \vspace{2pt}\textcolor{gold}{\rule{\linewidth}{1.2pt}}\par\smallskip
}
\newcommand{\worksmall}{\par\vspace{0.38in}}
\newcommand{\workmed}{\par\vspace{0.62in}}
\newcommand{\worklarge}{\par\vspace{0.92in}}
\newcommand{\workxl}{\par\vspace{1.22in}}
\newcommand{\choice}[1]{\(\square\)~#1\qquad}

\begin{document}
\summaryhead{8}{Convex Geometry}

\textbf{What this module is about.} Linear inequalities describe allowed regions rather than single lines. Each inequality selects one half-plane. When several constraints must all be satisfied, the feasible region is their intersection. In the two-variable problems in this module, feasible regions are usually convex polygons. Their vertices are especially important because a linear objective reaches its maximum or minimum at a corner point of a bounded feasible polygon.

\begin{keybox}{Convex sets and half-planes}
A set is \textbf{convex} if the complete line segment joining any two points in the set remains inside the set. Filled triangles, rectangles, circles, and half-planes are convex. Crescents, rings, and shapes with inward dents are not.

A linear inequality such as
\[
ax+by\le c
\]
has boundary line \(ax+by=c\) and one feasible half-plane. For \(\le\) or \(\ge\), the boundary is included. For \(<\) or \(>\), the boundary is not included.
\end{keybox}

\begin{examplebox}{Example 1: Find the boundary and test points}
Consider
\[
x+2y\le8.
\]
The boundary line \(x+2y=8\) has intercepts
\[
(8,0)\qquad\text{and}\qquad(0,4).
\]
Test the origin:
\[
0+2(0)=0\le8,
\]
so the feasible half-plane is the side containing \((0,0)\).

The point \((2,3)\) lies on the boundary because \(2+2(3)=8\). The point \((4,3)\) is not feasible because \(4+2(3)=10>8\).
\end{examplebox}

\begin{keybox}{Feasible regions and vertices}
To find a feasible polygon:
\begin{enumerate}[label=\arabic*.]
\item Graph or understand each boundary line.
\item Determine the feasible side of each inequality.
\item Intersect the allowed half-planes.
\item Find the vertices from axis intercepts and intersections of boundary lines.
\end{enumerate}
A proposed point is feasible only if it satisfies \emph{every} constraint.
\end{keybox}

\begin{examplebox}{Example 2: Finding a non-axis vertex}
Consider
\[
x+y\le8,
\qquad
2x+y\le10,
\qquad x\ge0,
\qquad y\ge0.
\]
The two non-axis boundary lines meet where
\[
x+y=8,
\qquad 2x+y=10.
\]
Subtracting the first equation from the second gives \(x=2\), and then \(y=6\). One vertex is therefore \((2,6)\). The complete feasible polygon has vertices
\[
(0,0),\ (5,0),\ (2,6),\ (0,8).
\]
\end{examplebox}

\begin{keybox}{Corner-point principle}
For a linear objective \(f=ax+by\) on a bounded feasible polygon, at least one maximum and at least one minimum occur at vertices. Therefore, evaluate the objective at every feasible vertex and compare the values.

If two \emph{adjacent} vertices have the same optimal value, then every point on the edge joining them is also optimal.
\end{keybox}

\prequizhead{8}{Convex Geometry}
\begin{refbox}{A reliable procedure}
Boundary line \(\to\) feasible side \(\to\) intersection of all constraints \(\to\) vertices \(\to\) evaluate the objective.
\end{refbox}

\begin{enumerate}
\item Consider a filled hexagon and a U-shaped region. Which one is convex? Use the definition of convexity to explain your answer in one sentence.
\workmed

\item A workshop constraint is
\[
2x+y\le10,
\]
where \(x\) and \(y\) are nonnegative quantities.
\begin{enumerate}[label=(\alph*),itemsep=4pt]
\item Find the positive \(x\)- and \(y\)-intercepts of the boundary line.
\item Determine whether \((3,4)\) is feasible.
\item Determine whether \((5,3)\) is feasible.
\end{enumerate}
\worklarge

\item Consider the system
\[
x+y\le9,
\qquad
2x+y\le12,
\qquad x\ge0,
\qquad y\ge0.
\]
Find the intersection of the two non-axis boundary lines. Then list all vertices of the feasible region.
\workxl

\item Using the feasible vertices from Problem 3,
\[
(0,0),\ (6,0),\ (3,6),\ (0,9),
\]
evaluate
\[
P=5x+4y
\]
at every vertex. State the maximum value and the vertex where it occurs.
\worklarge

\item A feasible polygon has vertices \((0,0),(5,0),(3,3),(0,6)\). Evaluate \(f=x+y\) at all four vertices. Two adjacent vertices tie for the maximum. What does that imply about the edge connecting them?
\worklarge
\end{enumerate}

\quizhead{8}{Convex Geometry}
\begin{enumerate}
\item Which is convex: a filled triangle or a ring-shaped region? Explain using line segments between points in the set.
\workmed

\item For the constraint \(3x+y\le12\), find the positive intercepts and determine whether \((3,4)\) is feasible.
\worklarge

\item Find all vertices of the feasible region
\[
x+y\le7,
\qquad
2x+y\le10,
\qquad x\ge0,
\qquad y\ge0.
\]
Show the calculation for the intersection of the two slanted boundary lines.
\workxl

\item Evaluate \(R=5x+4y\) at the vertices \((0,0),(5,0),(3,4),(0,7)\). Find the maximum and the vertex where it occurs.
\worklarge

\item Suppose a linear objective has the same maximum value at two adjacent vertices of a feasible polygon. Describe all of the optimal points.
\workmed
\end{enumerate}

\end{document}
